\documentclass[a4paper,11pt]{ltjsarticle}
\usepackage{luatexja}
\usepackage{luatexja-fontspec}
\usepackage{amsmath,amssymb,amsthm}
\usepackage{mathtools}
\usepackage{booktabs}
\usepackage{longtable}
\usepackage{array}
\usepackage{hyperref}
\usepackage{graphicx}
\usepackage{float}
\usepackage{geometry}
\geometry{margin=25mm}

\hypersetup{
  colorlinks=true,
  linkcolor=blue,
  urlcolor=blue,
  citecolor=blue
}

\theoremstyle{definition}
\newtheorem{theorem}{定理}[section]
\newtheorem{proposition}[theorem]{命題}
\newtheorem{lemma}[theorem]{補題}
\newtheorem{corollary}[theorem]{系}
\newtheorem{definition}[theorem]{定義}
\newtheorem{remark}[theorem]{注意}
\newtheorem{observation}[theorem]{観察}

\setlength{\parindent}{1em}
\setlength{\parskip}{0.5em}

\providecommand{\tightlist}{\setlength{\itemsep}{0pt}\setlength{\parskip}{0pt}}
\providecommand{\real}[1]{#1}
\begin{document}

\section{体積1の四次元超直方体に外接する超球体の直径は？}\label{ux4f53ux7a4d1ux306eux56dbux6b21ux5143ux8d85ux76f4ux65b9ux4f53ux306bux5916ux63a5ux3059ux308bux8d85ux7403ux4f53ux306eux76f4ux5f84ux306f}

\subsection{──
離散充填モデルから導かれる時空の次元・時間の起源・宇宙の始まり
──}\label{ux96e2ux6563ux5145ux586bux30e2ux30c7ux30ebux304bux3089ux5c0eux304bux308cux308bux6642ux7a7aux306eux6b21ux5143ux6642ux9593ux306eux8d77ux6e90ux5b87ux5b99ux306eux59cbux307eux308a}

\begin{center}\rule{0.5\linewidth}{0.5pt}\end{center}

\textbf{著者}: 木原 範昭 (WF System Co., Ltd.)

\textbf{日付}: 2026年4月（Rev.2）

\textbf{DOI：}
\href{https://doi.org/10.5281/zenodo.19834940}{10.5281/zenodo.19834940}（v2;
v1:
\href{https://doi.org/10.5281/zenodo.19533313}{10.5281/zenodo.19533313}）

\begin{center}\rule{0.5\linewidth}{0.5pt}\end{center}

\subsection{0.
本稿の主張範囲について（必読）}\label{ux672cux7a3fux306eux4e3bux5f35ux7bc4ux56f2ux306bux3064ux3044ux3066ux5fc5ux8aad}

本稿は\textbf{純粋に数学的な思考実験}であり、現実の宇宙について何かを主張するものではない。読者の誤読を避けるため、本稿が\textbf{主張すること}と\textbf{主張しないこと}を冒頭で明確にしておく。

\subsubsection{本稿が主張すること}\label{ux672cux7a3fux304cux4e3bux5f35ux3059ux308bux3053ux3068}

\begin{enumerate}
\def\labelenumi{\arabic{enumi}.}
\tightlist
\item
  「離散性」と「安定性」という2つの仮定、および「体積保存則」を出発点として、ある種の幾何学的・組合せ論的帰結が\textbf{演繹的に}導かれる、ということ。
\item
  その帰結の中に、\(\sqrt{n}\) が整数となる最小の非自明次元として
  \textbf{\(n = 4\)} が一意に選ばれる、という数学的事実が含まれること。
\item
  四次元における符号付き体積の反転操作には、\textbf{1軸反転}という最小経済的な選択肢が存在し、これが「1対3の非対称性」という構造を生むこと。
\item
  以上の数学的構造が、現実の四次元時空・ミンコフスキー計量・宇宙の始まりといった物理学上の概念と\textbf{構造的に類似した形を持つ}こと。
\item
  本稿の議論は本質的に\textbf{情報の組合せ論}であり、出発点となる「0次元状態」とは物理的な点ではなく、空間的自由度を持たずに静的に保持された\textbf{構造的情報}を指すこと（詳細は第8節の補足を参照）。
\end{enumerate}

\subsubsection{本稿が主張しないこと}\label{ux672cux7a3fux304cux4e3bux5f35ux3057ux306aux3044ux3053ux3068}

\begin{enumerate}
\def\labelenumi{\arabic{enumi}.}
\tightlist
\item
  \textbf{本稿は現実の宇宙が四次元である理由を説明するものではない。}
  現実の時空がなぜ4次元かは未解決の物理学的問題であり、本稿のモデルがその答えであるとは一切主張しない。
\item
  \textbf{本稿は物理学の理論ではない。}
  一般相対論・量子場理論・標準模型・弦理論・量子重力理論などとの厳密な対応は確立していない。本稿の「超直方体」が物理的実体に対応するという主張もない。
\item
  \textbf{本稿の「時間」「空間」「物質」「反物質」「ビッグバン」といった用語は比喩である。}
  物理学における同名の概念と同一視されるべきではない。あくまで数学的構造の特徴を説明するための借用語である。
\item
  \textbf{本稿は実証的予言を行わない。}
  観測・実験で検証可能な定量的予測を提供するものではなく、既存の物理理論に対する代替・反証・補強のいずれでもない。
\item
  \textbf{本稿は宇宙論的解釈を一切主張しない。}
  本稿に登場する宇宙論的な語り口（「宇宙の始まり」「時間の創発」「宇宙と反宇宙」等）は、すべて数学的構造の特徴を直観的に伝えるための\textbf{比喩的言語}であり、現実の宇宙に関する仮説・主張・予測ではない。
\end{enumerate}

\subsubsection{本稿の位置づけ}\label{ux672cux7a3fux306eux4f4dux7f6eux3065ux3051}

本稿は、四次元超直方体の外接球直径という素朴な幾何学的問いから出発し、対話的な思考実験を通じて、ごく少数の仮定からどこまで興味深い構造が導けるかを探究した記録である。\textbf{数学的遊戯}であり、\textbf{思考の体操}である。読者は本稿を、物理学的主張としてではなく、対称性・整数性・保存則といった概念がどのように相互作用するかを示す\textbf{例示}として読んでいただきたい。

数学的構造と物理的実在の対応を論じるには、本稿が扱う範囲をはるかに超えた厳密な定式化、既存物理学との接続、観測的検証が必要である。それらは本稿の射程外にある。

\subsubsection{改訂版（Rev.2）について}\label{ux6539ux8a02ux7248rev.2ux306bux3064ux3044ux3066}

初版§4は「一辺 \((2k+1)\)
の超立方体に外接する球」を扱っていたが、これは球の中に立方体形領域を置いただけで、球そのものを充填してはいなかった（充填率が常に
\(2/\pi^2 \approx 0.2026\) で頭打ちとなる）。改訂版では「半径
\(R = 2k+1\)
の球に単位立方体を稠密充填する」という本来の問いを定式化し直し、充填率が
\(k \to \infty\) で 1 に漸近収束する数列 \(N(k)\) を導出した。§7
以降の符号付き体積・時間軸・宇宙論的解釈の議論は、数列 \((2k+1)^4\)
を新しい充填数 \(N(k)\) に置き換える形でそのまま保持される。

\begin{center}\rule{0.5\linewidth}{0.5pt}\end{center}

\subsection{1.
素朴な問い：外接超球体の直径}\label{ux7d20ux6734ux306aux554fux3044ux5916ux63a5ux8d85ux7403ux4f53ux306eux76f4ux5f84}

一辺1の四次元超直方体（ハイパーキューブ）を考える。この超直方体に外接する四次元超球体、すなわちすべての頂点を通る超球体の直径はいくつか。

四次元超直方体の空間対角線の長さは

\[d = \sqrt{1^2 + 1^2 + 1^2 + 1^2} = \sqrt{4} = 2\]

である。外接超球体の直径はこの対角線に等しいため、\textbf{直径は2}（半径1）となる。

\begin{center}\rule{0.5\linewidth}{0.5pt}\end{center}

\subsection{2.
対称的な充填：何個詰め込めるか}\label{ux5bfeux79f0ux7684ux306aux5145ux586bux4f55ux500bux8a70ux3081ux8fbcux3081ux308bux304b}

次に、この四次元超直方体を複数用いて、回転対称性を維持したまま外接超球体の内部に詰め込むことを考える。

原点の周りに単位超立方体を対称に集めると、各超立方体が四次元空間の「象限」を一つずつ占めることができる。四次元では
\(2^4 = 16\)
個の象限が存在するため、\textbf{16個}の単位超立方体が原点で角を共有しながら隙間なく集まり、全体として一辺2の超立方体
\([-1,1]^4\) を構成する。

この配置は超立方体の対称群
\(B_4\)（位数384）を完全に保つ、最も対称的な詰め込みである。外接超球体の半径は中心から頂点までの距離
\(\sqrt{1^2+1^2+1^2+1^2} = 2\) であり、\textbf{直径は4}となる。

次元ごとの比較を示す：

{\def\LTcaptype{none} % do not increment counter
\begin{longtable}[]{@{}ccc@{}}
\toprule\noalign{}
次元 \(n\) & 個数 \(2^n\) & 外接球直径 \(2\sqrt{n}\) \\
\midrule\noalign{}
\endhead
\bottomrule\noalign{}
\endlastfoot
2 & 4 & \(2\sqrt{2} \approx 2.83\) \\
3 & 8 & \(2\sqrt{3} \approx 3.46\) \\
\textbf{4} & \textbf{16} & \textbf{4（整数）} \\
\end{longtable}
}

\begin{center}\rule{0.5\linewidth}{0.5pt}\end{center}

\subsection{3.
四次元の特殊性}\label{ux56dbux6b21ux5143ux306eux7279ux6b8aux6027}

ここで注目すべき事実が浮かび上がる。外接球の直径 \(2\sqrt{n}\)
が整数になるのは、\(\sqrt{n}\) が整数のとき、すなわち \(n\)
が完全平方数のときに限られる。

\[n = 1, 4, 9, 16, 25, \ldots\]

\(n = 1\) は自明な場合（1次元の線分）を除けば、\textbf{\(n = 4\)
が整数性を満たす最小の非自明な次元}である。これは偶然ではなく、以降の議論の基礎となる本質的な性質である。

\begin{center}\rule{0.5\linewidth}{0.5pt}\end{center}

\subsection{4.
稠密充填の定式化}\label{ux7a20ux5bc6ux5145ux586bux306eux5b9aux5f0fux5316}

\subsubsection{4.1 設定}\label{ux8a2dux5b9a}

原点中心の四次元球

\[B(R) := \{x \in \mathbb{R}^4 : \|x\|_2 \le R\}, \quad R = 2k+1 \;\; (k \in \mathbb{Z}_{\ge 0})\]

の内部に、中心が整数格子点 \((x, y, z, w) \in \mathbb{Z}^4\) にある一辺
1 の単位超立方体を、はみ出さないように詰める。各単位立方体は \(B(R)\)
に\textbf{完全に含まれる}ことを要請する。

\subsubsection{4.2 充填条件}\label{ux5145ux586bux6761ux4ef6}

中心 \((x, y, z, w)\) の単位立方体の頂点 16
点のうち、原点から最も遠いものは

\[P_{\max} = \left(|x|+\tfrac{1}{2},\ |y|+\tfrac{1}{2},\ |z|+\tfrac{1}{2},\ |w|+\tfrac{1}{2}\right)\]

である。立方体が \(B(R)\) に完全包含される必要十分条件は
\(\|P_{\max}\|_2 \le R\)、すなわち

\[\boxed{\;\left(|x|+\tfrac{1}{2}\right)^2 + \left(|y|+\tfrac{1}{2}\right)^2 + \left(|z|+\tfrac{1}{2}\right)^2 + \left(|w|+\tfrac{1}{2}\right)^2 \;\le\; (2k+1)^2\;}\]

この不等式を満たす整数組 \((x,y,z,w) \in \mathbb{Z}^4\) の個数を
\(N(k)\) と\textbf{定義}する。

\subsubsection{\texorpdfstring{4.3 数列
\(N(k)\)}{4.3 数列 N(k)}}\label{ux6570ux5217-nk}

数値計算により：

{\def\LTcaptype{none} % do not increment counter
\begin{longtable}[]{@{}cccc@{}}
\toprule\noalign{}
\(k\) & \(R = 2k+1\) & \(N(k)\) & 殻ごとの増分
\(a_k = N(k) - N(k-1)\) \\
\midrule\noalign{}
\endhead
\bottomrule\noalign{}
\endlastfoot
0 & 1 & 1 & --- \\
1 & 3 & 137 & 136 \\
2 & 5 & 1,545 & 1,408 \\
3 & 7 & 7,281 & 5,736 \\
4 & 9 & 22,409 & 15,128 \\
5 & 11 & 53,161 & 30,752 \\
6 & 13 & 108,081 & 54,920 \\
7 & 15 & 199,953 & 91,872 \\
8 & 17 & 337,417 & 137,464 \\
9 & 19 & 537,409 & 199,992 \\
10 & 21 & 818,145 & 280,736 \\
\end{longtable}
}

数列 \(N(k)\) は初版の四乗数列 \((2k+1)^4\)
ではなく、四次元球内の格子点計数に基づく数論的な数列である（後述§5）。

\subsubsection{4.4 充填率}\label{ux5145ux586bux7387}

四次元球の体積は \(V_4(R) = (\pi^2/2) R^4\)。充填率
\(\rho(k) := N(k) / V_4(2k+1)\) は：

{\def\LTcaptype{none} % do not increment counter
\begin{longtable}[]{@{}ccc@{}}
\toprule\noalign{}
\(k\) & \(V_4(R)\) & \(\rho(k)\) \\
\midrule\noalign{}
\endhead
\bottomrule\noalign{}
\endlastfoot
0 & 4.93 & 0.2026 \\
1 & 399.72 & 0.3427 \\
2 & 3,084.25 & 0.5009 \\
5 & 72,250.44 & 0.7358 \\
10 & 959,725.27 & 0.8525 \\
\end{longtable}
}

\textbf{命題 4.1}：\(\displaystyle \lim_{k \to \infty} \rho(k) = 1\).

\textbf{証明概略}：充填条件を満たす整数格子点の個数は、領域
\(\Omega_R := \{x \in \mathbb{R}^4 : \sum_{i=1}^{4} (|x_i|+1/2)^2 \le R^2\}\)
の体積に漸近する。\(\Omega_R\) は \(B(R)\)
を内側へ「半単位幅」だけ縮めた領域であり、その体積は \(R \to \infty\) で
\(V_4(R)\) に等しくなる。誤差項は表面積項 \(O(R^3)\)
で抑えられ、\(\rho(k) = 1 + O(R^{-1}) \to 1\)。\(\square\)

初版の充填率は \(k\) に依存せず定数 \(2/\pi^2 \approx 0.2026\)
だったが、改訂版では \(k\) の増加に伴い \(0.2026 \to 1\)
へ漸近的に成長する。これは「球が大きくなるほど隙間が単位立方体で埋められる」という直観に整合する。

\begin{center}\rule{0.5\linewidth}{0.5pt}\end{center}

\subsection{5.
四次元の特殊性（数論的背景）}\label{ux56dbux6b21ux5143ux306eux7279ux6b8aux6027ux6570ux8ad6ux7684ux80ccux666f}

\subsubsection{5.1
対角線解の整数性}\label{ux5bfeux89d2ux7ddaux89e3ux306eux6574ux6570ux6027}

四次元の対称解（対角線解）\(x = y = z = w\) を取ると

\[4x^2 = R^2 \quad\Longleftrightarrow\quad R = 2|x|\]

が成立する。これは \(\sqrt{4} = 2\)
が整数であることの直接的帰結であり、他の小次元では成立しない。

充填問題の文脈では、\(R = 2k+1\)（奇数）に対して半整数
\(|x| + 1/2 = (2k+1)/2\) を取ると

\[4 \cdot \left(\frac{2k+1}{2}\right)^2 = (2k+1)^2 = R^2\]

が成立する。

\subsubsection{5.2
角の立方体の球面接触}\label{ux89d2ux306eux7acbux65b9ux4f53ux306eux7403ux9762ux63a5ux89e6}

\textbf{命題 5.1}：各 \(k \ge 0\) について、中心
\((\pm k, \pm k, \pm k, \pm k)\) にある単位立方体（各象限に1個ずつ、計
\(2^4 = 16\) 個の「角の立方体」）は、最遠頂点が球面 \(\partial B(2k+1)\)
上に厳密に乗る。

\textbf{証明}：中心 \((\pm k)^4\) の立方体の最遠頂点は
\((\pm(k+1/2))^4\)。原点からの距離は

\[\sqrt{4 \cdot (k+1/2)^2} \;=\; 2(k+1/2) \;=\; 2k+1 \;=\; R \qquad \square\]

これは、四次元の単位立方体の頂点（半整数座標格子点）が外接球面上に整数距離で乗るという四次元固有の現象である。三次元では
\(\sqrt{3}\) が無理数のため成立しない。

\subsubsection{5.3 Lagrange--Jacobi
の数論的背景}\label{lagrangejacobi-ux306eux6570ux8ad6ux7684ux80ccux666f}

充填条件で変数変換 \(p_i := 2|x_i| + 1\)（奇正整数）を行うと、両辺に 4
を掛けて

\[p_1^2 + p_2^2 + p_3^2 + p_4^2 \;\le\; (4k+2)^2\]

が得られる。境界条件での等号成立 \(p_1^2 + \cdots + p_4^2 = (4k+2)^2\)
は、4 平方和の数論的問題に帰着する。

\textbf{Lagrange の四平方定理（1770）}：任意の非負整数は高々 4
個の平方の和で表せる。これは四次元が「全ての整数が四次元格子点で実現される」最小次元であることを意味する。

\textbf{Jacobi の四平方定理（1834）}：\(N\) を 4
個の整数の平方の和で表す方法の数 \(r_4(N)\) は閉形式

\[r_4(N) = \begin{cases} 8\,\sigma(N) & (4 \nmid N) \\ 24\,\sigma(N_{\text{odd}}) & (4 \mid N) \end{cases}\]

で与えられる（\(\sigma\) は約数和、\(N_{\text{odd}}\) は \(N\)
の奇成分）。この簡明な閉形式は四次元固有であり、三次元の \(r_3(N)\)
は類数を含むはるかに複雑な公式となる。

充填数 \(N(k)\) はこの数論的構造に直接連動する量である。

\begin{center}\rule{0.5\linewidth}{0.5pt}\end{center}

\subsection{6.
安定条件としての四次元（再定式化）}\label{ux5b89ux5b9aux6761ux4ef6ux3068ux3057ux3066ux306eux56dbux6b21ux5143ux518dux5b9aux5f0fux5316}

初版§6
の「外接超球体に対称的に充填でき、直径が整数となる」という安定条件を、改訂版では以下のように再定式化する。

\textbf{安定条件（改訂版）}：次元 \(n\) が以下をすべて満たすとき、\(n\)
は\textbf{安定}であると言う。

\begin{enumerate}
\def\labelenumi{\arabic{enumi}.}
\tightlist
\item
  半径 \(R \in \mathbb{Z}_{>0}\)
  の球への単位立方体稠密充填が定義できる。
\item
  充填率が \(R \to \infty\) で 1 に漸近収束する。
\item
  「角の立方体」の頂点が球面上の整数距離格子点上に乗る（\(\sqrt{n}\)
  の整数性）。
\item
  \(n\) 平方和の表現数 \(r_n(N)\) が単純な閉形式を持つ。
\end{enumerate}

これらすべてを満たす最小の非自明な次元は \textbf{\(n = 4\)} である。

\begin{itemize}
\tightlist
\item
  \(n = 3\) では条件 (3)(4) が成立しない。\(\sqrt{3}\)
  は無理数のため角の立方体は球面上に整数距離格子点として乗らず、\(r_3(N)\)
  も類数を含む複雑な公式となる。対称群 \(B_3\)（位数 48）も
  \(B_4\)（位数 384）と比べて貧弱である。
\item
  \(n = 9, 16, \ldots\) でも条件 (3) は成立するが、最小性で \(n = 4\)
  が選ばれる。
\end{itemize}

\begin{center}\rule{0.5\linewidth}{0.5pt}\end{center}

\subsection{7.
符号付き体積：マイナスの体積}\label{ux7b26ux53f7ux4ed8ux304dux4f53ux7a4dux30deux30a4ux30caux30b9ux306eux4f53ux7a4d}

ここで議論の方向を転換し、超立方体の体積に\textbf{符号}を導入する。

通常の幾何学的体積は常に非負であるが、線形代数的に符号付き体積（行列式で定義される向き付き体積）を考えることができる。四次元超立方体を4本のベクトル
\(\mathbf{v}_1, \mathbf{v}_2, \mathbf{v}_3, \mathbf{v}_4\)
で張られる平行体とみなすと、その符号付き体積は
\(\det[\mathbf{v}_1 \; \mathbf{v}_2 \; \mathbf{v}_3 \; \mathbf{v}_4]\)
である。

\(k\) 本のベクトルを反転（\(-1\) 倍）すると符号は \((-1)^k\)
倍される。したがって：

\begin{itemize}
\tightlist
\item
  1本反転：\((-1)^1 = -1\)（体積が負に）
\item
  2本反転：\((-1)^2 = +1\)（正のまま）
\item
  3本反転：\((-1)^3 = -1\)（負に）
\item
  4本反転：\((-1)^4 = +1\)（正に戻る）
\end{itemize}

負の体積は、1本または3本の軸を反転することで実現される。

\begin{center}\rule{0.5\linewidth}{0.5pt}\end{center}

\subsection{8.
初期条件：体積ゼロの圧縮状態}\label{ux521dux671fux6761ux4ef6ux4f53ux7a4dux30bcux30edux306eux5727ux7e2eux72b6ux614b}

ここから宇宙論的モデルの構築に入る。以下の初期条件を仮定する。

\textbf{宇宙の始まりが0次元から始まった}と仮定する。次元の自由度がない、極限まで圧縮された状態である。この状態に以下が詰め込まれている：

\begin{itemize}
\tightlist
\item
  体積 \(+1\) の四次元超直方体が \(N(k)\) 個
\item
  体積 \(-1\) の四次元超直方体が \(N(k)\) 個
\item
  \textbf{符号付き体積の合計はゼロ}
\end{itemize}

全体は半径 \(1\) の外接超球体の内側に圧縮されている。\(2 N(k)\)
個の超直方体が、本来 \(\rho(k)\) の充填率で詰まる体積 \(V_4(2k+1)\)
の球領域の情報を、自由度ゼロの一点に重畳した極限的な縮退状態である。

正負が打ち消し合って総体積ゼロでありながら、「個数」や「対称構造」としての情報は
\(N(k)\) という数に保持されている。

\subsubsection{補足：0次元における「情報」の意味について}\label{ux88dcux8db30ux6b21ux5143ux306bux304aux3051ux308bux60c5ux5831ux306eux610fux5473ux306bux3064ux3044ux3066}

ここで「0次元なのに \(N(k)\)
個の超直方体が存在するとはどういうことか」という疑問が当然生じる。これに対する本稿の立場を明確にしておく。

本稿における\textbf{0次元}とは、空間的な広がり（自由度）を一切持たない状態を指す。しかし0次元であることは、情報を持たないことを意味しない。0次元は「\textbf{ある}か\textbf{ない}か」という情報
\(I\) を担うことができる。

ここで言う「情報」とは、ビット列のような1と0の表現に限定されるものではない。実数値、複素数、行列、テンソル、あるいはそれらの組み合わせであってもよい。要件は単に：

\begin{quote}
\textbf{どのような情報が存在しているか／いないかが、空間的な広がりを伴わずに静的に表現されている}
\end{quote}

ということだけである。本稿の出発点における「\(N(k)\) 個の正の超直方体と
\(N(k)\)
個の負の超直方体」という記述は、空間に配置された物体を意味するのではなく、\textbf{この静的情報
\(I\)
が持つ構造}を記述したものである。個数・符号・対称性といった組合せ論的な情報が、自由度ゼロの状態として保持されている、と読んでいただきたい。

換言すれば、本稿の0次元状態は「点に物体が詰め込まれている物理的状況」ではなく、「\textbf{展開される前の構造的情報}」である。自由度の解放とは、この情報が次元を持った形で\textbf{読み出される}過程に対応する。

この意味で、本稿の議論は物理的描像である以前に\textbf{情報の組合せ論}であり、「0次元に情報が存在する」という言明は数学的に無矛盾である。

\begin{center}\rule{0.5\linewidth}{0.5pt}\end{center}

ここで最も重要な原理を明示する：

\begin{quote}
\textbf{体積保存則は破れない。}
\end{quote}

もし負の体積を許さなければ、宇宙の始まりは「体積ゼロから正の体積が湧き出す」ことになり、保存則が破られる。\textbf{負の体積は奇妙な仮定ではなく、保存則を守るための論理的必然}である。

\begin{center}\rule{0.5\linewidth}{0.5pt}\end{center}

\subsection{9.
最初の自由度：1次元への展開}\label{ux6700ux521dux306eux81eaux7531ux5ea61ux6b21ux5143ux3078ux306eux5c55ux958b}

次に、突然1次元の自由度が与えられたとする。このとき、以下の2つのルールを導入する：

\begin{enumerate}
\def\labelenumi{\arabic{enumi}.}
\tightlist
\item
  \textbf{充填率を最小化する}（できるだけ薄く広がる）
\item
  \textbf{四次元超体積の合計を最小化する}（四次元方向の広がりを最小にする）
\end{enumerate}

両方のルールを同時に満たすと、正負のペアが1次元軸上に点状に並ぶのが極限的な解となる。各ペアの四次元断面は限りなくゼロに近づき、長さ方向にのみ
\(N(k)\) だけ伸びる。

結果として全体は「太さゼロ・長さ \(N(k)\)
の1次元的な線」に縮退する。四次元の自由度を捨てて1次元の自由度に乗り換える転移であり、元の四次元構造の情報は長さという1次元量
\(N(k)\) に完全にエンコードされて残る。

\begin{center}\rule{0.5\linewidth}{0.5pt}\end{center}

\subsection{10.
2次元への展開と直径最小化}\label{ux6b21ux5143ux3078ux306eux5c55ux958bux3068ux76f4ux5f84ux6700ux5c0fux5316}

自由度が2次元になった場合、3つ目のルールを追加する：

\begin{enumerate}
\def\labelenumi{\arabic{enumi}.}
\setcounter{enumi}{2}
\tightlist
\item
  \textbf{外接超球体の直径を最小化する}
\end{enumerate}

1次元のままでは長さ \(N(k)\)
がそのまま直径となり、巨大すぎる。2次元の自由度があれば、これを\textbf{正方形状に折りたたむ}ことで直径を最小化できる。

\(N(k)\) 個を正方形に並べると一辺は

\[\sqrt{N(k)} = N(k)^{1/2}\]

となり、外接超球の直径は正方形の対角線

\[\sqrt{2} \cdot N(k)^{1/2}\]

に縮まる。四次元方向は依然として点状に潰れたまま、2次元方向に
\(N(k)^{1/2} \times N(k)^{1/2}\) の格子状に展開された状態である。

\begin{center}\rule{0.5\linewidth}{0.5pt}\end{center}

\subsection{11.
次元ごとの一辺の長さのパターン}\label{ux6b21ux5143ux3054ux3068ux306eux4e00ux8fbaux306eux9577ux3055ux306eux30d1ux30bfux30fcux30f3}

ここでパターンが見えてくる。\(N(k)\) 個を \(n\)
次元方向に展開すると、一辺の長さは \(N(k)^{1/n}\)
となる。\(N(k) \sim (\pi^2/2) R^4\) (\(R = 2k+1\), \(k \to \infty\))
を用いると、漸近形は次のようになる：

{\def\LTcaptype{none} % do not increment counter
\begin{longtable}[]{@{}ccc@{}}
\toprule\noalign{}
自由度の次元 \(n\) & 一辺の長さ & 漸近形 (\(k \to \infty\)) \\
\midrule\noalign{}
\endhead
\bottomrule\noalign{}
\endlastfoot
1 & \(N(k)\) & \((\pi^2/2)\,(2k+1)^4\) \\
2 & \(N(k)^{1/2}\) & \((\pi^2/2)^{1/2}\,(2k+1)^2\) \\
3 & \(N(k)^{1/3}\) & \((\pi^2/2)^{1/3}\,(2k+1)^{4/3}\) \\
\textbf{4} & \(N(k)^{1/4}\) & \((\pi^2/2)^{1/4}\,(2k+1)\) \\
\end{longtable}
}

\(n = 4\) で漸近形が球の半径 \(R = 2k+1\) に\textbf{比例}する（係数
\((\pi^2/2)^{1/4} \approx 1.4894\)）。これは単位立方体格子のスケールが球の半径と漸近的に一致することを意味し、対称的な四次元超立方体構造の\textbf{漸近的回復}として再解釈される。\(n = 3\)
では漸近形が \((2k+1)^{4/3}\)
となり半径より速く発散して次元の整合性が破れる。\(n < 4\)
ではいずれも一辺と半径のスケールが整合しない。

\begin{center}\rule{0.5\linewidth}{0.5pt}\end{center}

\subsection{12.
なぜ四次元に落ち着くのか}\label{ux306aux305cux56dbux6b21ux5143ux306bux843dux3061ux7740ux304fux306eux304b}

ルール群（対称性の維持、\(\sqrt{n}\)
の整数性による角の立方体の球面接触、外接球直径の最小化、超体積の最小化）をすべて満たす「落ち着ける次元」を探すと：

\begin{itemize}
\tightlist
\item
  \textbf{\(n = 0\)}：始まりの状態。総体積ゼロ、半径 1
  の球に圧縮。出発点。
\item
  \textbf{\(n = 1\)}：長さ \(N(k)\)
  の線。直径が爆発的に大きくなり、直径最小化に反する。通過点。
\item
  \textbf{\(n = 2\)}：\(\sqrt{n}\)
  が無理数で角の立方体が球面に厳密に内接できない中間状態。
\item
  \textbf{\(n = 3\)}：\(\sqrt{n}\) が無理数。一辺 \(N(k)^{1/3}\)
  は球の半径 \(R\) と整合的なスケールにならない。不安定。
\item
  \textbf{\(n = 4\)}：\(\sqrt{4} = 2\)
  で角の立方体が球面に厳密に内接（命題 5.1）。\(B_4\)
  対称性完全回復。一辺は漸近的に \((\pi^2/2)^{1/4}(2k+1)\)
  に収束し、球の半径 \(R = 2k+1\) と同じスケール。直径 \(2(2k+1)\)
  で安定。\textbf{最初の真の平衡点。}
\end{itemize}

\(n = 9, 16, \ldots\)
にもさらなる平衡点は存在するが、直径最小化と次元を上げないバランスにおいて、最小の解である
4 で止まる。

\begin{quote}
\textbf{0次元から始まり自由度を順に解放していくと、1〜3次元はどれも安定して留まれず、四次元に到達して初めてすべてのルールが同時に満たされて落ち着く。}
\end{quote}

この系列を見ると、0の次に落ち着ける次元は4であり、自由度を与えると結果として四次元で安定する。四次元であることは仮定ではなく帰結である。

\begin{center}\rule{0.5\linewidth}{0.5pt}\end{center}

\subsection{13.
時間軸の起源：なぜ1軸だけが特別か}\label{ux6642ux9593ux8ef8ux306eux8d77ux6e90ux306aux305c1ux8ef8ux3060ux3051ux304cux7279ux5225ux304b}

負の体積を実現するためには、奇数本の軸を反転する必要がある。最も経済的な操作は\textbf{1本だけの反転}であり、\((-1)^1 = -1\)
を生む。

4つの軸のうち\textbf{ちょうど1本だけ}を他と区別して反転することで負体積が実現される。残り3本は反転に関与せず対等なまま。

\begin{quote}
\textbf{この「1対3の非対称性」が、四次元時空における時間軸1本と空間軸3本という構造の起源である。}
\end{quote}

ミンコフスキー計量 \(ds^2 = -dt^2 + dx^2 + dy^2 + dz^2\)
において時間成分だけが負の符号を持つ。本モデルでは、この符号の非対称性が「負体積を作るために1軸だけ反転した」という起源論的事実から\textbf{自然に導かれる}。計量に後から符号を入れるのではなく、符号の非対称性が宇宙の始まりの構造そのものに由来する。

三次元立方体モデルではこの一意性は成立しない。三次元では
\((-1)^1 = -1\), \((-1)^3 = -1\)
であり、1本反転と3本反転のどちらも負を生じる。「1対2」と「3対0」のどちらが選ばれるか決まらず、時間軸の特別性が一意に定まらない。四次元の偶数性が、1本だけ反転
= 時間軸が一意に決まるという構造を保証している。

\begin{center}\rule{0.5\linewidth}{0.5pt}\end{center}

\subsection{14.
体積保存則と宇宙の始まり}\label{ux4f53ux7a4dux4fddux5b58ux5247ux3068ux5b87ux5b99ux306eux59cbux307eux308a}

モデルの論理構造を改めて整理する：

\begin{enumerate}
\def\labelenumi{\arabic{enumi}.}
\tightlist
\item
  体積保存則は破れない（物理の大原則）
\item
  → 始まりがあるなら総体積は常にゼロでなければならない
\item
  → 負の体積が論理的に要請される
\item
  → 軸反転が必要
\item
  → 偶数次元でなければペア対称性が保てない
\item
  → 安定には \(\sqrt{n}\)
  の整数性により角の立方体が球面上の格子点に乗ることが要請される
\item
  → 最小の非自明な解は \(n = 4\)
\item
  → 負体積の最小操作は1軸反転
\item
  → 1軸だけが特別 = 時間
\item
  → 残り3軸が対等 = 空間
\end{enumerate}

\textbf{仮定は2つだけである：}

\begin{itemize}
\tightlist
\item
  \textbf{仮定1}：宇宙は離散的である
\item
  \textbf{仮定2}：安定な最小単位が存在する
\end{itemize}

四次元であること、時間軸の特殊性、宇宙の始まりは、いずれもこの2つの仮定からの\textbf{帰結}であり、入力パラメータではない。

\begin{center}\rule{0.5\linewidth}{0.5pt}\end{center}

\subsection{15.
球殻モデル：正曲率の時空}\label{ux7403ux6bbbux30e2ux30c7ux30ebux6b63ux66f2ux7387ux306eux6642ux7a7a}

ここでルールを一つ変更する。「充填率を下げる（内部に均一に広がる）」ではなく、「次元の端に集まる」というルールにするとどうなるか。

四次元超球の表面は三次元球面 \(S^3\)
である。超直方体の正負ペアが内部に均一に広がるのではなく表面に集中すると、\(N(k)\)
個のペアが半径 \((2k+1)\) の超球の表面に極薄の殻として張り付く。

安定条件はすべて満たされる：

\begin{itemize}
\tightlist
\item
  \(\sqrt{4} = 2\) の整数性は維持
\item
  \(B_4\) 対称性も球殻上で維持
\item
  外接球直径 \(2(2k+1)\) も不変
\item
  正負ペアの総体積ゼロも保存
\end{itemize}

この変形の決定的に重要な点は、\textbf{我々が観測する三次元空間の起源}を自然に説明できることである。四次元超球の表面は三次元多様体であるため、この薄い殻の上に住む観測者にとっては宇宙は三次元空間に見える。

この描像は\textbf{正の曲率を持つ閉じた四次元時空}のメタファーとなっている。一般相対論における閉じたFriedmann宇宙（\(k = +1\)
のFLRW計量）では、各時刻における空間断面が \(S^3\)
であり、その半径が時間とともに変化する。本モデルでは：

\begin{itemize}
\tightlist
\item
  \(S^3\) の半径 \(= (2k+1)\) → 時間発展に対応
\item
  殻の厚み \(\approx 0\) → 物質が空間断面上に集中
\item
  正の曲率 \(= 1/(2k+1)^2\) → \(k\) が大きくなると曲率が減少
\end{itemize}

\(k\) が増大すると曲率 \(1/(2k+1)^2\)
が小さくなり、局所的にはほぼ平坦に見える。これは現在の宇宙が観測上ほぼ平坦であるにもかかわらず、厳密にはわずかに正の曲率を持つ可能性があるという観測事実と整合する。

\begin{center}\rule{0.5\linewidth}{0.5pt}\end{center}

\subsection{16.
宇宙と反宇宙：正負の分離}\label{ux5b87ux5b99ux3068ux53cdux5b87ux5b99ux6b63ux8ca0ux306eux5206ux96e2}

さらにルールを導入する。

\begin{quote}
\textbf{外接超球体は複数あってもよく、全体で体積が保存されればよい。}
\end{quote}

この制約の緩和により、正と負の超直方体を同じ超球に詰め込む必要がなくなる。すると：

\begin{itemize}
\tightlist
\item
  \textbf{正の超球}：直径 \(2(2k+1)\)、\(N(k)\) 個の体積 \(+1\)
  の超直方体
\item
  \textbf{負の超球}：直径 \(2(2k+1)\)、\(N(k)\) 個の体積 \(-1\)
  の超直方体
\item
  \textbf{総体積}：\(+N(k) - N(k) = 0\)（保存則を満たす）
\end{itemize}

両方とも同じ直径・同じ個数・同じ対称性を持ち、それぞれ独立に安定条件を満たす。

これは\textbf{宇宙と反宇宙のペア生成}を意味する。0次元状態からの展開が、正負が完全に分離した双子の宇宙として実現される。

\begin{center}\rule{0.5\linewidth}{0.5pt}\end{center}

\subsection{17.
時間を持つ宇宙と時間を持たない宇宙}\label{ux6642ux9593ux3092ux6301ux3064ux5b87ux5b99ux3068ux6642ux9593ux3092ux6301ux305fux306aux3044ux5b87ux5b99}

1軸反転が時間を生むという先の議論に立ち返る。

\textbf{我々の宇宙（負の体積側）}： - 1軸が反転 → \((-1)^1 = -1\) →
負の体積 - 時間軸が存在する → 動的で変化がある - ミンコフスキー計量
\((-,+,+,+)\)

\textbf{ペアの宇宙（正の体積側）}： - 軸の反転がない → \((+1)^4 = +1\) →
正の体積 - 時間軸がない → 4軸すべてが対等 - ユークリッド計量
\((+,+,+,+)\)

ペアの宇宙は時間を持たないため、変化も因果も動力学もない。四次元の空間として静的に「在る」だけの構造である。我々から観測することは原理的にできない。観測という行為自体が時間を前提とするからである。

\begin{quote}
\textbf{時間の存在そのものが「負債」であり、それを相殺する「時間なき静的構造」が必ず対になって存在しなければならない。}
\end{quote}

時間は「あって当たり前のもの」ではなく、保存則を満たすための代償として生まれた構造である。

\begin{center}\rule{0.5\linewidth}{0.5pt}\end{center}

\subsection{18.
初期配分の揺らぎと宇宙の複雑さ}\label{ux521dux671fux914dux5206ux306eux63faux3089ux304eux3068ux5b87ux5b99ux306eux8907ux96d1ux3055}

2つの超球への正負の配分が完全に対称でない初期条件を考える。

\textbf{正の超球に} \(+N\) 個と \(-m\) 個、\textbf{負の超球に} \(-N\)
個と \(+m\)
個が混在している場合（ただし総和は常にゼロ）、各超球内部で「少数派の反対符号」が存在する不安定な状態となる。

この少数派が多数派と対消滅していく過程で、以下の現象が自然に出現する：

\textbf{対消滅のダイナミクス}：我々の超球（負の体積側）に混入した正の体積超直方体が、負のものと対になって消えていく。これは初期宇宙での物質・反物質の対消滅に対応する。

\textbf{非対称性の残存}：配分が完全対称でなく微小な偏りがあれば、対消滅後にわずかな「残り」が生じる。これがバリオン非対称性（現在の宇宙に物質だけが残っている理由）に対応する。観測では約10億個に1個の割合で物質が反物質より多かった。本モデルでは初期配分の微小なずれとして表現される。

\textbf{構造形成の種}：配分のずれのパターンによっては、局所的に正負が分離しきれず塊を形成する。これが密度ゆらぎ、ひいては銀河や大規模構造の種となりうる。

\textbf{複雑さのパラメータ}：\(k\)
が小さければ超直方体の総数が少なく比較的早く安定する。\(k\)
が大きければ組み合わせが膨大になり、複雑で豊かな中間過程が生まれる。\(k\)
は宇宙の「複雑さ」を決めるパラメータとなる。

\begin{center}\rule{0.5\linewidth}{0.5pt}\end{center}

\subsection{19.
時間の創発：最初は時間がなかった}\label{ux6642ux9593ux306eux5275ux767aux6700ux521dux306fux6642ux9593ux304cux306aux304bux3063ux305f}

最後に、最も深い帰結を述べる。

負の体積を実現する方法は1軸反転だけではない。3軸反転でも \((-1)^3 = -1\)
で負となる。初期条件で負体積の超直方体の反転軸がバラバラだった場合を考える。

\textbf{初期状態}：ある超直方体は第1軸を反転、別のものは第3軸を反転、また別のものは3軸同時反転。「どの軸が時間か」が定まっていない。どの方向を見ても対等であり、\textbf{一貫した時間軸が存在しない。つまり時間がない。}

\textbf{中間過程}：対消滅と再配置が進むにつれ、最もエネルギー的に経済的な配置（=1軸反転のみ）が優勢になっていく。これは\textbf{対称性の自発的破れ}である。4軸が対等だった状態から1本だけが選ばれていく。

\textbf{安定状態}：ほぼすべての超直方体が同じ1軸を反転軸として共有するようになる。この瞬間、初めて一貫した時間方向が宇宙全体に定義される。\textbf{時間の創発}である。

\begin{quote}
\textbf{時間は宇宙の基本構造ではなく、安定化の結果として創発した秩序である。}
\end{quote}

初期には時間がなく四次元空間として対等だった。対消滅のダイナミクスを通じて反転軸が揃い、時間が凝縮した。

この描像は以下の物理学的問題と接点を持つ：

\begin{itemize}
\tightlist
\item
  \textbf{プランク時代}：初期宇宙では時空の概念そのものが成立しないとされる。本モデルでは反転軸が揃っていないため時間が定義できない状態として自然に表現される。
\item
  \textbf{インフレーション}：反転軸が急速に揃っていく過程が、指数関数的膨張に対応する可能性がある。「時間が生まれる」ことと「空間が急膨張する」ことが同じ過程の二つの側面となる。
\item
  \textbf{量子重力との接点}：多くの量子重力理論で「時間は基本的ではなく創発的」とされている。本モデルはこれを離散的な超直方体の反転軸の統計的整列として具体的に実現する。
\end{itemize}

\begin{center}\rule{0.5\linewidth}{0.5pt}\end{center}

\subsection{20.
結論：2つの仮定から導かれるもの}\label{ux7d50ux8ad62ux3064ux306eux4eeeux5b9aux304bux3089ux5c0eux304bux308cux308bux3082ux306e}

本モデルの仮定を厳密に数えると：

\begin{itemize}
\tightlist
\item
  \textbf{仮定1}：宇宙は離散的である
\item
  \textbf{仮定2}：安定な最小単位が存在する
\end{itemize}

ここから体積保存則を経由して以下がすべて\textbf{演繹的に導かれる}：

\begin{enumerate}
\def\labelenumi{\arabic{enumi}.}
\tightlist
\item
  \textbf{時空が四次元であること} --- \(\sqrt{n}\)
  が整数になる最小の非自明次元
\item
  \textbf{時間軸が1本だけ特別であること} ---
  負体積の最小操作としての1軸反転
\item
  \textbf{空間が3次元で等方的であること} --- 反転に関与しない残りの3軸
\item
  \textbf{ミンコフスキー計量の符号構造} \((-,+,+,+)\) --- 反転軸の符号
\item
  \textbf{宇宙が「始まれた」こと} --- 体積ゼロからの保存則を破らない展開
\item
  \textbf{宇宙と反宇宙のペア構造} --- 正負分離による安定化
\item
  \textbf{物質・反物質の非対称性} --- 初期配分の微小なずれ
\item
  \textbf{時間の創発} --- 反転軸の統計的整列
\item
  \textbf{正曲率の閉じた宇宙の可能性} --- 球殻モデルとの対応
\end{enumerate}

次元の数すら導出量であって入力パラメータではない。

本モデルは、追加の物理定数や余剰次元のコンパクト化や人間原理を必要としない。離散性と安定性という2つの仮定と、体積保存則という1つの原理から、時空の構造と宇宙の起源に関する主要な問いに統一的に答える枠組みを提供する。

\begin{center}\rule{0.5\linewidth}{0.5pt}\end{center}

\emph{本稿は、四次元超直方体の外接球直径という素朴な幾何学的問いから出発し、対話的な思考実験を通じて段階的に構築されたモデルの記録である。}

\begin{center}\rule{0.5\linewidth}{0.5pt}\end{center}

\subsection{付録:
シリーズの論理的読み順}\label{ux4ed8ux9332-ux30b7ux30eaux30fcux30baux306eux8ad6ux7406ux7684ux8aadux307fux9806}

本シリーズの論文は公開順に番号付けされているが、論理的な読み順としては以下を推奨する:

\begin{enumerate}
\def\labelenumi{\arabic{enumi}.}
\tightlist
\item
  論文 {[}1{]}--{[}3{]}:
  基礎フレームワーク（中心投影、対称性、観測の相対性）
\item
  論文 {[}4{]}: \(R\) 極限と入れ子構造
\item
  論文 {[}5{]}: 次元追加
\item
  論文 {[}9{]}: 論文 {[}1{]} の方程式の Schwarzschild--de Sitter
  厳密解と \(R \to 0\) の地平面構造
\item
  論文 {[}7{]}: 測地線の橋渡し
\item
  論文 {[}10{]}: \(R \to 0\)
  極限の次元的解釈（入れ子構造・次元階層・厳密解の統合）
\item
  論文 {[}6{]}: 各次元の主観空間の記述
\item
  \textbf{本稿（論文 {[}8{]}）}: 0次元からの次元展開
\end{enumerate}

\begin{center}\rule{0.5\linewidth}{0.5pt}\end{center}

\end{document}
